\section{计算图上的反向传播与自动微分}
\label{sec:10-Matrix-Calculus/04-Applications/03}

上一节的两层网络你已经手推过一遍，代码也跑通了。接着把隐藏层从两层加到二十四层，接上残差和 LayerNorm，同一份脚本报 OOM。batch 砍掉一半才勉强跑得动，而显存监控给出的分布不太符合直觉：占大头的不是参数，是前向留下的一堆中间张量。有人让你打开梯度检查点，显存立刻降下来一多半，每步时间涨了约三成。

三个现象挤在一起：反向为什么要用前向的中间量，这些中间量为什么比参数还占地方，重算又凭什么能把显存换回来。它们指向同一个结构。\hyperref[sec:10-Matrix-Calculus/02-Differentials/02]{``链式法则就是线性映射复合''}已经算清了那笔总账——复合的导数是矩阵连乘，连乘的括号可以任意打，打在输出端那一侧代价低一个维度。这一节把那条结论安置到具体的图上：节点是算子，边是局部导数，反向传播就是这串连乘按从右往左的顺序求值。

先划一条界。\hyperref[sec:08-Numerical-Analysis/07-Numerical-Methods-for-ML/02]{``自动微分给出导数，梯度检查建立信任''}处理的是自动微分与数值差分的关系：为什么差商既慢又不准，一次梯度检查要怎样设计才有判断力。这里不重复那一侧。这一节只回答一个问题：本 Part 前面几节手推出来的矩阵形式的梯度，怎样变成图上一条可以机械执行、且与图的全局结构无关的规则。

\begin{center}
\begin{tikzpicture}[x=1cm,y=1cm,font=\small,>=stealth]
  \node[font=\footnotesize,black!60,anchor=west] at (-0.6,3.05) {前向：从左往右求值，\(X,Z_1,H\) 留在显存里等反向来取};
  \foreach \x/\lab in {0/{\(X\)},2.3/{\(Z_1\)},4.6/{\(H\)},6.9/{\(Z_2\)}}
    {\draw[draw=black!50,fill=black!5] (\x,2.0) circle (0.42);
     \node[font=\footnotesize] at (\x,2.0) {\lab};}
  \draw[draw=black!50,fill=orange!18] (9.2,2.0) circle (0.42);
  \node[font=\footnotesize] at (9.2,2.0) {\(L\)};
  \draw[draw=black!55,fill=blue!8,rounded corners=2pt] (1.85,0.2) rectangle (2.75,0.9);
  \node[font=\footnotesize] at (2.3,0.55) {\(W_1\)};
  \draw[draw=black!55,fill=blue!8,rounded corners=2pt] (6.45,0.2) rectangle (7.35,0.9);
  \node[font=\footnotesize] at (6.9,0.55) {\(W_2\)};
  \foreach \a/\b in {0/2.3,2.3/4.6,4.6/6.9,6.9/9.2}
    \draw[->,black!55] (\a+0.42,2.0) -- (\b-0.42,2.0);
  \draw[->,black!55] (2.3,0.9) -- (2.3,1.55);
  \draw[->,black!55] (6.9,0.9) -- (6.9,1.55);
  \node[font=\footnotesize,black!70] at (1.15,2.35) {仿射};
  \node[font=\footnotesize,black!70] at (3.45,2.35) {\(\phi\)};
  \node[font=\footnotesize,black!70] at (5.75,2.35) {仿射};
  \node[font=\footnotesize,black!70] at (8.05,2.35) {交叉熵};
  \node[font=\footnotesize,red!55!black] at (1.15,1.58) {\(\cdot\,W_1^{\trans}\)};
  \node[font=\footnotesize,red!55!black] at (3.45,1.58) {\(\odot\,\phi'(Z_1)\)};
  \node[font=\footnotesize,red!55!black] at (5.75,1.58) {\(\cdot\,W_2^{\trans}\)};
  \node[font=\footnotesize,red!55!black] at (8.05,1.58) {\(\tfrac1n(P-Y)\)};
  \draw[->,red!60!black,line width=1.1pt] (9.2,-0.6) -- (0,-0.6);
  \node[font=\footnotesize,red!55!black,anchor=south] at (4.6,-0.5) {反向：从右往左，每一步都是一次``行向量乘矩阵''，\(O(d^2)\)};
  \node[font=\footnotesize,red!55!black,anchor=north] at (2.3,-0.78) {\(\bar W_1=X^{\trans}\bar Z_1\)};
  \node[font=\footnotesize,red!55!black,anchor=north] at (6.9,-0.78) {\(\bar W_2=H^{\trans}\bar Z_2\)};
\end{tikzpicture}

\emph{每个节点只认得红色那一格，也就是自己的局部导数；反向沿相反的顺序把它们乘起来，中间结果始终与激活同形，完整的 Jacobian 一次也没有出现。}
\end{center}

图里没有一个箭头知道整张网络长什么样。这正是这套机制能扩展到几百个算子的原因：算子的作者只需要交出局部导数，顺序由图负责。下面把顺序、开销和显存三件事分别算清楚。

\subsection{从输出端起步，中间结果才始终是一行}

把网络看成 \(L\) 个映射的复合，第 \(k\) 层的 Jacobian 记为 \(J_k\)，为了估开销先假设它们都是 \(d\times d\)。损失对输入的 Jacobian 是连乘 \(J=J_L\cdots J_1\)。

从输入端起步意味着先算 \(J_2J_1\)，再乘 \(J_3\)，如此下去。每一次都是矩阵乘矩阵，\(O(d^3)\)，而且中途必须腾出 \(O(d^2)\) 的空间放那个方阵。从输出端起步则不同。损失是标量，最右端交出来的是一个行向量——图里那个 \(\bar Z_2\) 摊平就是它——于是每一步都是行向量乘矩阵，\(O(d^2)\)，中间结果永远只有 \(d\) 个数。

差的这一个 \(d\) 倍不是常数因子上的胜负。\(d\) 取几千时，从输入端起步的做法连一层都算不完。所以顺序不是实现细节，它就是反向模式与前向模式的分界：同一串连乘，从右往左求值叫反向模式，传播的是 VJP；从左往右求值叫前向模式，传播的是 JVP。两者结果完全相同，因为矩阵乘法结合律不挑边。

具体到图上，反向模式的执行规则短到可以背下来。给每个节点配一个伴随量，形状与该节点的值一致；损失节点的伴随播下种子 \(\bar L=1\)；每个节点从下游收到自己的伴随，乘上自己那个局部 Jacobian 的转置，把结果交给上游；有多条下游边流入的节点，把各路收到的伴随相加。走完拓扑逆序，所有参数节点手上的就是梯度。

规则里没有出现``网络''这个词。节点不需要知道自己身处第几层、后面接的是什么、损失是哪一种。这是分工上的关键一点：写一个新算子时你要交的只是它的前向和它的局部 VJP，剩下的由图的遍历完成。上一节那个 softmax 与交叉熵的融合算子，正是这个接口下的一次手工优化——把两个节点合成一个，局部 VJP 从 Jacobian 相乘变成 \(P-Y\)。

\subsection{两层网络的反向，逐节点写一遍}

把图里的箭头换成式子。设 \(X\in\R^{n\times d}\)、\(W_1\in\R^{d\times h}\)、\(b_1\in\R^h\)、\(W_2\in\R^{h\times K}\)、\(b_2\in\R^K\)，前向是

\[
\begin{aligned}
Z_1&=XW_1+\mathbf1b_1^{\trans}, &\quad H&=\phi(Z_1),\\
Z_2&=HW_2+\mathbf1b_2^{\trans}, &\quad L&=\tfrac1n\textstyle\sum\nolimits_i\ell(Z_2[i,:],Y[i,:]).
\end{aligned}
\]

沿用加横线的记号表示损失对某个量的梯度。融合的 softmax 交叉熵节点直接交出上一节的结果 \(\bar Z_2=\frac1n(P-Y)\)，反向从这里往左走：

\[
\begin{aligned}
\bar W_2&=H^{\trans}\bar Z_2, &\quad \bar b_2&=\bar Z_2^{\trans}\mathbf1, &\quad \bar H&=\bar Z_2W_2^{\trans},\\
\bar Z_1&=\bar H\odot\phi'(Z_1), &&&&\\
\bar W_1&=X^{\trans}\bar Z_1, &\quad \bar b_1&=\bar Z_1^{\trans}\mathbf1, &\quad \bar X&=\bar Z_1W_1^{\trans}.
\end{aligned}
\]

七个式子，三条规则。矩阵乘法 \(AB\) 的反向对 \(A\) 乘右因子的转置、对 \(B\) 乘左因子的转置；逐元素函数的反向仍是逐元素，对角 Jacobian 退化成 \(\odot\)，代价 \(O(nh)\) 而不是 \(O(n h^2)\)；广播前向复制，反向就沿被复制的那个轴求和，\(\mathbf1b^{\trans}\) 的反向是 \(\bar Z^{\trans}\mathbf1\) 正是这一条。整个 Part 里所有的手推，落到图上就只剩这三条。

推导本身不长，容易出错的是条件。先对形状：\(\bar W_2\) 应当是 \(h\times K\)，\(H^{\trans}\bar Z_2\) 恰好是 \((h\times n)(n\times K)\)；\(\bar H\) 应当是 \(n\times h\)，只有 \(\bar Z_2W_2^{\trans}\) 摆得进去。形状是必要条件而不是充分条件——\(h=K\) 时几种错误摆法同样能通过，真正定住答案的是上一单元的迹整理。再对缩放：\(1/n\) 只在损失那一个节点上出现一次，若它同时被写进了 \(\bar W_1\) 的表达式，梯度会小 \(n\) 倍，而损失曲线看上去仍然在降。然后对求和轴：bias 的梯度是沿样本轴求和，不是求平均，平均已经由 \(\bar Z_2\) 里的 \(1/n\) 承担了。最后对逐元素项：\(\phi'\) 必须在 \(Z_1\) 处取值，不是在 \(H\) 处——除非你用的是 \(\tanh\) 或 sigmoid 这种导数能写成输出函数的激活，那时框架会存输出而丢掉输入，这是一个算子级的省显存手法，也是下面那一小节的伏笔。

\subsection{输入维数远小于输出维数时，前向模式更省}

反向模式在深度学习里的统治地位来自一个特定的形状：参数上亿，损失是一个标量。方向数少的那一端播种，标量损失这一端只有一个方向，所以一趟走完。反过来的场景不是假想的。

考虑对少数几个标量做灵敏度分析：一个温度系数、一个正则强度、一个物理常数，输出却是几万个网格点或整条轨迹上的值。此时输入方向只有一个，输出方向上万，前向模式跑一趟就能拿到所有输出对这个输入的导数，反向模式要跑上万趟。前向模式在图上的走法与反向恰好镜像：给输入节点播下一个切向量，沿边的正方向传播，每个节点把上游的切向量乘上自己的局部 Jacobian 交给下游。

这个镜像还带来一个常被忽略的好处。前向模式与前向计算同步进行，一个节点用完就可以丢，不需要为反向保留任何中间量。它的显存开销是常数倍的前向开销，没有下面那笔账。

两种模式还能叠起来用。Hessian 与向量的乘积可以写成对梯度再做一次前向传播，前向套反向各一趟，代价与一次普通反向同阶，全程不构造那个 \(d\times d\) 的矩阵，细节见\hyperref[sec:10-Matrix-Calculus/02-Differentials/03]{``Hessian 是二阶线性算子''}。判断规则始终是同一条：数一数你要的导数涉及几个输入方向、几个输出方向，从少的那一端播种。

\subsection{显存的账记在``反向要用前向的量''这一条上}

现在可以回答开头那个 OOM 了。看一眼上面那七个式子里出现了哪些前向量：\(\bar W_2\) 要 \(H\)，\(\bar H\) 要 \(W_2\)，\(\bar Z_1\) 要 \(Z_1\)，\(\bar W_1\) 要 \(X\)。乘法的反向需要另一个因子，逐元素函数的反向需要它的输入或输出——局部 VJP 几乎都不是纯函数，它依赖前向在那一点算出来的值。所以前向必须把这些中间量留在显存里，一直留到反向走到它为止。

这笔开销的形状与参数完全不同。参数量固定，中间量却随 batch 和序列长度线性增长。一个宽度 \(4096\)、batch 与序列长度乘积为三万的层，一个激活张量就是一亿多个数，半精度下约 \(0.25\) GB；二十四层每层留一个，就是六个 GB，而且这还只算了每层一个。batch 一加大先爆的是它，不是参数，开头那个显存分布就是这么来的。

梯度检查点做的是拿计算换存储。只在若干个位置存下激活，其余的丢掉；反向走到某一段时，从最近的那个检查点重新前向一遍，把段内的中间量重算出来，用完再丢。\(L\) 层网络每隔 \(\sqrt L\) 层设一个检查点，显存从 \(O(L)\) 降到 \(O(\sqrt L)\)，代价是多做一次前向的计算量。前向大约占一次训练步的三分之一，所以时间涨三成上下——正是开头看到的数字。

同一条取舍上还有几个别的位置。逐元素激活存输出不存输入，能省掉一半；算子融合让一段连续运算只在边界处落盘；可逆结构从下一层的输出反解出上一层的输入，把存储彻底换成计算。它们改的都是同一个式子里的两项配比，而不是别的东西。

反过来，任何让前向中间量失效的操作都会在这里出事。原地修改是最常见的一种：把 \(H\) 覆盖掉再进入反向，\(\bar W_2=H^{\trans}\bar Z_2\) 里用到的就是被改过的值。框架通常会检测到并报错，但也存在检测不到的写法，尤其是自定义算子里手动管理缓冲区的时候。

\subsection{图上求导，求的是实际执行的那个程序的导数}

最后收一下边界。自动微分保证的是``对这一次真正跑过的那串运算求导''，这句话有两处需要留意。

一处是运算本身。\(\relu\) 在零点、\(\max\) 在并列点、取整与排序都没有唯一的导数，框架返回的是一个约定值；含随机采样的模型每次前向对应的函数都不同，反向拿到的是这一次那个函数的梯度而不是期望的梯度；\texttt{stop\_gradient} 直接命令链式法则在此截断。这几类情形在\hyperref[sec:10-Matrix-Calculus/02-Differentials/02]{``链式法则就是线性映射复合''}里已经讨论过，结论是它们都不算实现错误，只是``执行的程序''与``你想求导的函数''不再是同一个。

另一处是图的构造。控制流让每一次的图形状都可能不同，框架必须记录实际走过的分支；截断的迭代求解器有两种反向语义，对执行的有限步求导与对收敛方程隐式求导给出不同的结果，两者都合法但不相等，选哪一种是建模决定而非实现细节。

所以手推没有被自动微分免除，只是换了用途。它不再用来生成公式，而用来在框架给出一个数时判断这个数对不对：形状对不对、缩放对不对、融合有没有把中间量省掉、显存的账该记在哪一行。本 Part 从``导数是最佳线性映射''走到这里，工具已经交齐；\hyperref[sec:14-Deep-Learning/02-Losses-and-Backpropagation/02]{``一次完整的前向与反向传播''}会在一个完整网络上把这套流程从头到尾再走一遍，那时每一步的来历都应当是认得出的。
